\section{Uvod u proširenu stvarnost}

Proširena stvarnost (engl. \textit{Augmented Reality}, skraćeno AR) predstavlja tehnologiju koja digitalni sadržaj --- trodimenzionalne modele, tekst, slike ili animacije --- projektuje preko slike stvarnog sveta u realnom vremenu. Za razliku od klasičnih računarskih prikaza, koji postoje isključivo na ekranu i odvojeno od fizičkog okruženja, proširena stvarnost digitalne elemente uklapa u prostor oko korisnika tako da oni deluju kao njegov sastavni deo: virtuelna fotelja stoji na stvarnom podu, digitalna strelica pokazuje duž stvarne ulice, a informaciona tabla lebdi iznad stvarnih vrata učionice. Ključna odlika ove tehnologije jeste interaktivnost u realnom vremenu --- kako se korisnik kreće kroz prostor, digitalni sadržaj se kontinuirano prilagođava novoj perspektivi, čime se stvara utisak da virtuelni objekti zaista postoje u fizičkom svetu.

Proširenu stvarnost treba jasno razlikovati od virtuelne stvarnosti (engl. \textit{Virtual Reality}). Virtuelna stvarnost korisnika u potpunosti izmešta u sintetičko okruženje: stavljanjem namenskih naočara stvarni svet nestaje, a zamenjuje ga računarski generisana scena. Proširena stvarnost, nasuprot tome, stvarni svet zadržava kao osnovu i samo ga dopunjuje. Ovaj odnos slikovito opisuje kontinuum stvarnosti i virtuelnosti --- zamišljena osa na čijem jednom kraju stoji potpuno stvarno okruženje, a na drugom potpuno virtuelno. Proširena stvarnost nalazi se bliže stvarnom kraju kontinuuma, jer virtuelni elementi čine manji deo doživljaja, dok virtuelna stvarnost zauzima suprotni pol. Između njih se prostiru različiti oblici mešane stvarnosti, u kojima se udeo digitalnog i fizičkog sadržaja postepeno menja.

Iako se proširena stvarnost doživljava kao savremena tehnologija, njeni koreni sežu decenijama unazad. Još krajem šezdesetih godina dvadesetog veka konstruisani su prvi uređaji koji su se nosili na glavi i pred oči korisnika postavljali jednostavnu računarsku grafiku uklopljenu u pogled na prostoriju. Ti rani sistemi bili su glomazni, vezani kablom za računar i dostupni isključivo u laboratorijama. Tokom devedesetih godina tehnologija je našla primenu u vojnoj avijaciji i industrijskim istraživanjima, a sam termin „proširena stvarnost" ušao je u širu upotrebu. Prava prekretnica, međutim, nastupila je pojavom pametnih telefona: uređaj koji milijarde ljudi svakodnevno nose u džepu poseduje upravo ono što je proširenoj stvarnosti potrebno --- kameru, ekran, senzore pokreta i dovoljno snažan procesor. Tehnologija koja je nekada zahtevala specijalizovanu i skupu opremu postala je dostupna svakom vlasniku prosečnog mobilnog telefona.

Široka dostupnost otvorila je vrata izuzetno raznovrsnim primenama. U oblasti igara i zabave globalni fenomen postala je igra \textit{Pokémon GO}, koja je virtuelna bića smestila u parkove, ulice i trgove stvarnih gradova i time milionima korisnika po prvi put približila ideju proširene stvarnosti. U trgovini i opremanju doma aplikacije poput \textit{IKEA Place} omogućavaju da se komad nameštaja u prirodnoj veličini „postavi" u dnevnu sobu pre kupovine, pa kupac odmah vidi da li se nova sofa uklapa u prostor i bojom i dimenzijama. U navigaciji se uputstva za kretanje iscrtavaju direktno preko slike ulice koju kamera snima, tako da pešak umesto apstraktne mape vidi strelicu koja lebdi tačno iznad raskrsnice na kojoj treba da skrene.

Obrazovanje je još jedno područje u kojem proširena stvarnost pokazuje punu vrednost: umesto statičnih ilustracija u udžbeniku, učenici mogu da obiđu trodimenzionalni model ljudskog srca koji kuca na školskoj klupi ili da posmatraju Sunčev sistem kako kruži iznad katedre. U medicini se tehnologija koristi za vizuelizaciju vena prilikom davanja injekcija, kao i za navođenje hirurga tokom zahvata, pri čemu se snimci unutrašnjih organa projektuju direktno preko tela pacijenta. U industriji i održavanju, tehničar koji popravlja složenu mašinu može kroz ekran ili pametne naočare da vidi uputstva za sklapanje iscrtana tačno preko delova na koje se odnose, čime se smanjuje broj grešaka i skraćuje obuka.

Pored ovih specijalizovanih domena, proširena stvarnost nezametno poboljšava i svakodnevni život. Kamera telefona prevodi natpise i jelovnike na stranom jeziku i prikazuje prevod na samoj slici, a aplikacije za probu naočara ili šminke pokazuju kako će proizvod izgledati na licu korisnika. Posebno zanimljivu primenu, međutim, otvaraju pametna okruženja: savremene zgrade, kampusi i učionice opremljeni su senzorima koji neprekidno mere temperaturu, buku i kvalitet vazduha, ali ti podaci po pravilu ostaju skriveni u udaljenim kontrolnim tablama i izveštajima. Proširena stvarnost nudi daleko prirodniji pristup tim informacijama --- dovoljno je uperiti telefon u vrata učionice da bi se iznad njih, u stvarnom prostoru, prikazali živi podaci o toj prostoriji: da li je zauzeta i do kada, koji se čas u njoj trenutno održava i kakvi su uslovi za rad. Upravo ta primena --- prikaz kontekstualnih informacija iz pametnog okruženja u proširenoj stvarnosti --- predmet je aplikacije opisane u ovom radu.

Razlog zbog kojeg je proširena stvarnost na pametnim telefonima doživela tako brzu ekspanziju jeste ukidanje ulaznih barijera. Korisniku nije potreban nijedan dodatni uređaj: dovoljan je telefon koji već poseduje, a programerima su na raspolaganju zrele softverske platforme koje složene probleme računarskog vida --- praćenje položaja uređaja, prepoznavanje površina i procenu osvetljenja --- rešavaju umesto njih. Na \textit{Android} platformi centralno mesto među tim rešenjima zauzima \textit{ARCore}, tehnologija kompanije \textit{Google} na kojoj je zasnovana i aplikacija iz ovog rada, a koja je detaljno predstavljena u narednom poglavlju.
